\subsection*{把比例增长推广成向量之间的耦合}
标量方程 $x'=ax$ 用指数描述。若状态包含多个分量，每个分量的速度都由全部当前分量的线性组合决定，就得到 $x'=Ax$。需要寻找一个随时间变化的矩阵，把初值送到后续状态，矩阵指数承担这个任务。

把标量指数的幂级数中数字 $a$ 换成矩阵 $A$，得到 $\mathrm e^{tA}$。它不是对每个元素单独取指数，因为矩阵乘法描述分量间反复耦合，$A^2,A^3$ 保留了这种传递结构。

\subsection*{从能截断的例子建立信心}
取 $A=\left[\begin{smallmatrix}0&1\\0&0\end{smallmatrix}\right]$，对应方程 $x'=y,y'=0$。第二分量固定，第一分量以固定速度移动，所以解应是 $(x_0+ty_0,y_0)$。由于 $A^2=0$，指数级数恰好截断为 $I+tA$，给出同一结果。

先演化 $s$ 再演化 $t$，等于演化 $s+t$。在本例中展开乘积，额外的 $stA^2$ 为零，便直接验证了这个性质。它表达的是同一个自治系统连续运行的时间相加。

但不同矩阵代表不同耦合规则，通常不能把 $\mathrm e^{A+B}$ 拆成 $\mathrm e^A\mathrm e^B$，除非有相应交换条件。与矩阵乘法一样，动作顺序不可随意忽略。这个简单系统是入口，一般矩阵指数仍需要级数收敛和求导的分析基础。
